\documentclass[11pt,a4paper]{report}
\usepackage[margin=2.2cm]{geometry}
\usepackage{booktabs,tabularx,longtable}
\usepackage[table]{xcolor}
\usepackage{enumitem,fancyhdr,hyperref}
\newcolumntype{P}[1]{>{\raggedright\arraybackslash}p{#1}}
\definecolor{ink}{HTML}{1E3348}\definecolor{green}{HTML}{247A67}\definecolor{pale}{HTML}{EEF6F3}
\hypersetup{colorlinks=true,linkcolor=green,urlcolor=green}
\pagestyle{fancy}\fancyhf{}\lhead{Northstar Foods | Quarterly statement}\rhead{\thepage}
\renewcommand{\headrulewidth}{0.4pt}
\setlength{\headheight}{14pt}
\begin{document}
\begin{titlepage}\centering\vspace*{2cm}{\color{ink}\Huge\bfseries NORTHSTAR FOODS\par}\vspace{0.5cm}{\Large Corporate quarterly financial statement\par}\vspace{1.7cm}{\color{green}\fontsize{32}{38}\selectfont Q1 2026\par}\vspace{1cm}\large Management reporting pack\\Issued 15 April 2026
\vfill\begin{tabular}{rl}Reporting basis:& Management accounts\\Currency:& EUR millions\\Prepared by:& Finance planning and analysis\end{tabular}
\vfill{\small Fictional example data. This starter is not audited financial advice.}\end{titlepage}
\chapter{Quarter at a glance}
Revenue grew 6.3\% year on year, led by the chilled range. Operating margin was below plan because of a temporary packaging surcharge and launch costs. Cash remains adequate, but the next quarter depends on converting the March order book without extending customer terms.
\section{Key indicators}
\begin{center}\rowcolors{2}{pale}{}
\begin{tabular}{lrrr}\toprule Indicator & Q1 2025 & Q1 2026 & Plan variance\\\midrule Revenue & 5.8 & 6.2 & +0.1\\ Gross margin & 34.1\% & 33.7\% & -0.8pp\\ Operating margin & 12.0\% & 11.4\% & -1.1pp\\ Closing cash & 2.7 & 2.4 & -0.2\\\bottomrule\end{tabular}\end{center}
\chapter{Financial statements}
\section{Income statement}
\begin{center}\rowcolors{2}{pale}{}
\begin{tabular}{lrrr}\toprule EUR m & Q1 2025 & Q1 2026 & Change\\\midrule Revenue & 5.80 & 6.17 & 6.3\%\\ Cost of goods sold & (3.82) & (4.09) & 7.1\%\\ Gross profit & 1.98 & 2.08 & 5.1\%\\ Operating expenses & (1.28) & (1.38) & 7.8\%\\ Operating profit & 0.70 & 0.70 & 0.0\%\\\bottomrule\end{tabular}\end{center}
\section{Balance-sheet snapshot}
\begin{center}\begin{tabular}{lrr}\toprule EUR m & 31 Dec 2025 & 31 Mar 2026\\\midrule Cash and equivalents & 2.8 & 2.4\\ Trade receivables & 2.1 & 2.3\\ Inventory & 3.4 & 3.8\\ Property and equipment & 6.2 & 6.3\\ Trade payables & (2.0) & (2.2)\\ Borrowings & (4.1) & (4.0)\\\bottomrule\end{tabular}\end{center}
\chapter{Cash bridge and commentary}
\begin{center}\rowcolors{2}{pale}{}
\begin{tabular}{lr}\toprule EUR m & Q1 2026\\\midrule Opening cash & 2.8\\ Operating profit & 0.7\\ Working-capital investment & (0.6)\\ Capital expenditure & (0.4)\\ Financing and other & (0.1)\\ \textbf{Closing cash} & \textbf{2.4}\\\bottomrule\end{tabular}\end{center}
\section{Variance notes}
\begin{enumerate}[leftmargin=*]\item Packaging costs were EUR 0.09m above plan; a supplier reset is scheduled for May.\item Receivables increased with two large March shipments; collection dates are confirmed.\item Capital expenditure is on schedule for the line-efficiency project and remains within the approved envelope.\end{enumerate}
\chapter{Outlook and controls}
The Q2 forecast assumes 7\% volume growth, stable pricing, and a return to 48-day receivables. Management will review the order book, cash forecast, and headroom against the EUR 1.2m minimum liquidity threshold every Friday.
\begin{longtable}{@{}P{0.29\textwidth}P{0.31\textwidth}P{0.27\textwidth}@{}}\toprule Control & Owner & Next review\\\midrule
Forecast refresh & FP\&A lead & Weekly\\
Receivables ageing & Controller & 22 April\\
Capex gate & Operations director & 30 April\\
Board outlook & CFO & 6 May\\\bottomrule\end{longtable}
\end{document}
